% chapters/ch04-conceptual-to-relational.tex
\chapter{From Conceptual Model to Relational Schema}

\begin{learningobjectives}
In this chapter you will learn how to translate conceptual models into relational schemas in a disciplined way.
You will be able to map entities and relationships to relations, choose foreign key placement correctly, and recognise when you need association relations.
You will also learn how to write a schema specification that is meaningful even before SQL is introduced.
\end{learningobjectives}

\section{Translation as rule-based work}
Once the domain is clear and identity is defined, moving to a relational schema should not feel like invention.
It should feel like applying rules.
This mindset matters because beginners often “design by improvisation”, adding columns and tables until the model seems to work.
That approach produces schemas that are hard to justify and even harder to maintain.

A relational schema is a statement about:
what exists (relations), what uniquely identifies it (primary keys), how it connects (foreign keys), and what is allowed (constraints).
The translation step is where meaning becomes enforceable structure.

\subsection{Schema notation used in this book}
We will specify relations without SQL, using a compact notation.
For example:
\[
\begin{aligned}
\rel{Appointment}(&\underline{appointment\_id}, start\_time, \\
&patient\_id \rightarrow \rel{Patient}, \\
&doctor\_id \rightarrow \rel{Doctor})
\end{aligned}
\]
Underline indicates the primary key.
An arrow indicates a foreign key reference.
Optional attributes may be written with a question mark (e.g., \texttt{return\_date?}) when the idea of optionality matters.

This notation is intentionally simple.
It is designed for clarity and discussion.
It forces you to state identity and references explicitly.

\section{The core mapping rules}
The following rules cover most designs you will meet at introductory level.
They are written as design rules, not as SQL steps.

\begin{principlebox}
\textbf{Core mapping rules}
\begin{enumerate}[label=\arabic*.]
  \item Each entity becomes a relation.
  \item Each relation has a primary key that represents identity.
  \item A 1:N relationship becomes a foreign key on the N side.
  \item A 1:1 relationship becomes a foreign key on one side, with a uniqueness constraint.
  \item An N:M relationship becomes an association relation with two foreign keys.
  \item Multivalued attributes become separate relations.
  \item Attributes of a relationship belong in the structure that represents that relationship.
\end{enumerate}
\end{principlebox}

The final rule is often the most important.
If a property describes the association, not either entity alone, then it belongs with the relationship.
Dates on a loan are not properties of the member or the copy; they are properties of the borrowing event.

\section{Mapping entities and relationships}

\subsection{Mapping entities: the easy part}
Mapping an entity is straightforward.
You create a relation and choose its primary key.
The harder part is deciding which attributes are actually attributes, and which are hidden entities.

For example, in a clinic story, “specialty” might look like an attribute of Doctor.
That can be fine for a minimal model.
But if specialties have their own identity, descriptions, and governance (and doctors can have multiple specialties), then Specialty is an entity, not a string attribute.
Translation works best when those conceptual decisions are made explicitly.

\subsection{Mapping 1:N relationships: foreign keys and meaning}
A one-to-many relationship means that one instance on the “one” side can be related to many on the “many” side.
In the relational schema, this is represented by storing the primary key of the “one” side as a foreign key in the “many” side relation.

The important point is not the mechanical placement of the foreign key.
The important point is that the foreign key expresses a constraint: each “many” row points to exactly one “one” row (if the relationship is mandatory).

\begin{examplebox}
\textbf{Clinic: Patient \(\rightarrow\) Appointment is 1:N}\par\smallskip
Conceptual meaning: a patient can have many appointments; each appointment has one patient.\par\smallskip
Relational schema:
\[
\rel{Patient}(\underline{patient\_id}, name, birth\_date)
\]
\[
\rel{Appointment}(\underline{appointment\_id}, start\_time, patient\_id \rightarrow \rel{Patient})
\]
\end{examplebox}

Optionality matters here.
If an appointment must have a patient, then \texttt{patient\_id} is mandatory.
If an appointment could exist without a patient (usually not true), then \texttt{patient\_id} would be optional.
This is not a technical detail.
It is a statement about what data states are valid.

\subsection{Mapping 1:1 relationships: choose a direction, then enforce it}
One-to-one relationships require care.
A 1:1 relationship means that each instance on either side is associated with at most one instance on the other side.
In the relational schema, you can represent this with a foreign key on one side, but you must enforce uniqueness to maintain the one-to-one meaning.

Where you place the foreign key is a design choice guided by optionality and lifecycle.
If one side can exist without the other, place the foreign key on the optional side.
This avoids NULLs on the mandatory side and usually matches creation workflows.

\begin{examplebox}
\textbf{Example: Patient \(\leftrightarrow\) InsuranceProfile (1:1)}\par\smallskip
If a patient may have no insurance profile yet, then:
\[
\rel{Patient}(\underline{patient\_id}, \dots)
\]
\[
\begin{aligned}
\rel{InsuranceProfile}(&\underline{insurance\_profile\_id}, \\
&patient\_id \rightarrow \rel{Patient} \;[\text{unique}], \\
&provider, policy\_number)
\end{aligned}
\]
Uniqueness on \texttt{patient\_id} in \rel{InsuranceProfile} enforces the 1:1 relationship.
\end{examplebox}

The main mistake is to store both foreign keys in both tables.
That duplicates the relationship and creates inconsistency risk.
A one-to-one relationship should be represented once, with the right constraint.

\begin{principlebox}
\textbf{One-to-one placement rule}\\
Place the foreign key on the optional side and enforce uniqueness there.
This preserves the 1:1 meaning without introducing mandatory NULLs.
\end{principlebox}

\subsection{Mapping N:M relationships: association relations}
Many-to-many relationships cannot be represented as a single foreign key without losing information.
The relational model represents them using an association relation (also called a junction table).
This relation contains at least two foreign keys, one to each participating entity.
Its primary key is often the pair of foreign keys.

\begin{examplebox}
\textbf{Doctor \(\leftrightarrow\) Specialty (N:M)}\par\smallskip
\[
\rel{Doctor}(\underline{doctor\_id}, name)
\]
\[
\rel{Specialty}(\underline{specialty\_id}, name)
\]
\[
\begin{aligned}
\rel{DoctorSpecialty}(&\underline{doctor\_id} \rightarrow \rel{Doctor}, \\
&\underline{specialty\_id} \rightarrow \rel{Specialty})
\end{aligned}
\]
\end{examplebox}

This structure becomes even more important when the relationship has attributes.
For instance, if a doctor has a specialty \emph{since a date}, that date belongs to the association relation:
\[
\rel{DoctorSpecialty}(\underline{doctor\_id}, \underline{specialty\_id}, since\_date)
\]
This follows the principle: relationship attributes belong where the relationship is represented.

\begin{notebox}
If you add relationship attributes to an entity instead of the association relation, you silently change the meaning.
Dates and roles are almost never properties of a single side.
\end{notebox}

\subsection{Multivalued attributes are a hidden relationship}
When you see an attribute that can take multiple values for a single entity, you are not looking at a single attribute.
You are looking at a separate structure.
In relational design, this becomes its own relation, typically with a foreign key back to the owning entity.

For example, if a member can have multiple phone numbers, “phone\_numbers” is not an attribute in the relational sense.
A separate relation is appropriate:
\[
\rel{MemberPhone}(\underline{member\_id} \rightarrow \rel{Member}, \underline{phone\_number})
\]
Whether \texttt{phone\_number} is part of the primary key depends on the domain rules, but the modelling decision is the key point.

\section{Relationship attributes and event entities}
Some relationships are really events.
If a relationship has time, status, or quantities, it usually deserves its own entity.
This keeps the attributes where they belong and preserves history.

\begin{examplebox}
\textbf{Loan as event, not just a FK}\par\smallskip
Instead of storing a \texttt{current\_copy\_id} on \rel{Member}, represent the borrowing as:
\[
\begin{aligned}
\rel{Loan}(&\underline{loan\_id}, checkout\_date, due\_date, return\_date?, \\
&member\_id \rightarrow \rel{Member}, copy\_id \rightarrow \rel{Copy})
\end{aligned}
\]
This captures the event and allows multiple loans over time.
\end{examplebox}

\begin{figure}[htbp]
\centering
\includegraphics[width=0.85\textwidth]{relationship-event.pdf}
\caption{When a relationship has its own attributes, treat it as an event entity.}
\end{figure}

\section{Optionality mapping rules}
Optionality should be decided explicitly.
It changes where you place foreign keys and whether nulls are acceptable.

Rule A: If the relationship is mandatory on the “many” side, the FK is mandatory.\\
Rule B: If the relationship is optional, the FK may be optional.\\
Rule C: In 1:1, place the FK on the optional side.

\begin{examplebox}
\textbf{Optionality fragments}\par\smallskip
Mandatory 1:N:
\[
\rel{Appointment}(\underline{appointment\_id}, patient\_id \rightarrow \rel{Patient})
\]
Optional 1:N:
\[
\rel{Reservation}(\underline{reservation\_id}, member\_id? \rightarrow \rel{Member})
\]
Optional 1:1:
\[
\rel{InsuranceProfile}(\underline{insurance\_profile\_id}, patient\_id \rightarrow \rel{Patient} \;[\text{unique}])
\]
\end{examplebox}

\begin{figure}[htbp]
\centering
\includegraphics[width=0.85\textwidth]{optionality-mapping.pdf}
\caption{Three common ways to represent optional information in a relational schema.}
\end{figure}

\section{A worked example: Library loans and reservations}
The library domain is useful because it forces you to distinguish between a \term{Title} and a physical \term{Copy}.
It also forces you to model events over time, such as loans, which are relationship-like structures with their own attributes.

Start with entities:
Member, Title, Copy.
A Title has many Copies (1:N).
A Member borrows a Copy over time, which is best represented as a Loan entity because it carries dates and status.

The conceptual meaning “member borrows copy” is not a static relationship.
It is a time-indexed event that must be recorded.
That is why \rel{Loan} exists.

\begin{examplebox}
\textbf{Library schema (logical, no SQL)}\par\smallskip
\[
\rel{Member}(\underline{member\_id}, name, email \;[\text{unique}])
\]
\[
\rel{Title}(\underline{title\_id}, isbn \;[\text{unique}], title, author)
\]
\[
\begin{aligned}
\rel{Copy}(&\underline{copy\_id}, barcode \;[\text{unique}], status, \\
&title\_id \rightarrow \rel{Title})
\end{aligned}
\]
\[
\begin{aligned}
\rel{Loan}(&\underline{loan\_id}, checkout\_date, due\_date, return\_date?, \\
&member\_id \rightarrow \rel{Member}, copy\_id \rightarrow \rel{Copy})
\end{aligned}
\]
\[
\begin{aligned}
\rel{Reservation}(&\underline{reservation\_id}, request\_date, status, \\
&member\_id \rightarrow \rel{Member}, title\_id \rightarrow \rel{Title})
\end{aligned}
\]
\end{examplebox}

The important constraints are not cosmetic.
They preserve meaning:
a Copy belongs to exactly one Title,
a Loan references exactly one Member and one Copy,
and a Copy should not have two active loans at the same time.
That last rule is not automatically guaranteed by foreign keys; it is a business constraint you must state and later enforce.

\section{Mini-case: E-commerce orders}
An order is an event with its own identity.
Line items are not attributes of the order; they are entities that link products to orders with quantities.

\begin{examplebox}
\textbf{Order and order line (conceptual)}\par\smallskip
\[
\rel{Order}(\underline{order\_id}, placed\_at, customer\_id \rightarrow \rel{Customer})
\]
\[
\begin{aligned}
\rel{OrderLine}(&\underline{order\_id} \rightarrow \rel{Order}, \\
&\underline{product\_id} \rightarrow \rel{Product}, quantity, unit\_price)
\end{aligned}
\]
The line is the association; quantities belong there, not on \rel{Order}.
\end{examplebox}

\section{Common mistakes}
There are recurring errors at this stage.

Mistake: Collapsing Title and Copy into a single relation.\\
Fix: Separate entities so availability and history can be represented cleanly.

Mistake: Treating time-indexed events as static relationships.\\
Fix: Model the event as its own relation (Loan, Appointment, Enrollment).

Mistake: Forgetting optionality.\\
Fix: Decide mandatory vs optional and place FKs accordingly.

Mistake: Putting relationship attributes on an entity.\\
Fix: Move dates, quantities, and roles to the association relation.

\section{Checklist: mapping rules you can apply mechanically}
Use these rules as a quick mechanical pass:
\begin{itemize}[leftmargin=*,nosep]
  \item Each entity becomes one relation.
  \item Each relation has a primary key that represents identity.
  \item 1:N becomes a foreign key on the N side.
  \item 1:1 becomes a foreign key with a uniqueness constraint.
  \item N:M becomes an association relation with two FKs.
  \item Multivalued attributes become their own relations.
\end{itemize}

\section{Mini checklist}
Before finalizing the relational schema, check:
\begin{itemize}[leftmargin=*,nosep]
  \item Does every relationship map to exactly one structure in the schema?
  \item Are all 1:1 relationships enforced with uniqueness, not duplication?
  \item Are relationship attributes stored on the association relation?
  \item Are multivalued attributes separated into their own relations?
  \item Are mandatory and optional references explicit?
  \item Can you explain why each foreign key is placed where it is?
  \item Are event entities used for time- or status-rich relationships?
\end{itemize}

\begin{chaptersummary}
Translating a conceptual model into a relational schema is a disciplined application of mapping rules.
Entities become relations with keys.
One-to-many relationships become foreign keys on the many side.
Many-to-many relationships become association relations.
If a relationship has attributes or time semantics, it usually needs its own structure.
A schema is correct when it preserves meaning through identity, references, and constraints, even before SQL is introduced.
\end{chaptersummary}

\begin{exercisebox}
Choose a domain (clinic, library, university, online shop) and answer in full sentences.
\begin{enumerate}[label=\arabic*.]
  \item Identify one 1:N relationship and write the corresponding two relations in book notation.
  \item Identify one N:M relationship and write the association relation.
  \item Identify one relationship that has attributes (e.g., dates) and explain where those attributes belong.
  \item Identify one multivalued attribute and rewrite it as a separate relation.
  \item Write three business constraints that are not automatically enforced by foreign keys.
  \item Show a 1:1 relationship and explain where you placed the foreign key and why.
  \item Give one example where optionality changes the schema and explain the difference.
  \item Sketch an order/order-line schema and justify where quantities and prices belong.
  \item Write a constraint that enforces “one active loan per copy” in words.
\end{enumerate}
\end{exercisebox}

In the next chapter we examine redundancy and normalization.
The mapping choices here determine where redundancy appears and how it can be removed.
